\begin{bild}{Integrierte Rechnung v1.6: von Pauli zu Weyl, aber noch nicht zu Raum}
Die Compiler-Matrizen liefern drei hermitesche $\alpha_j$ mit $\{\alpha_j,\alpha_k\}=2\delta_{jk}I$. Daraus entstehen exakt die acht Koeffizienten
\[
A_\epsilon=\frac18(I+\epsilon_1\alpha_1)(I+\epsilon_2\alpha_2)(I+\epsilon_3\alpha_3).
\]
Ergänzt man räumliche Verschiebungen, erhält man einen unitären Walk und in erster Ordnung einen Weyl-Operator. Ohne diese Ergänzung ist $\sum_\epsilon A_\epsilon=I$. Die Matrixbrücke verkleinert damit die fehlende Stelle: Nicht weitere interne Pauli-Matrizen, sondern die physisch erzeugten Verschiebungen, ihre gemeinsame Skala und die Beseitigung der zusätzlichen Moden sind zu begründen. Die exakte Rechnung steht in Kapitel~\ref{ch:v16-operations}.
\end{bild}
